\section{Hardware Implementation and Support}
\label{sec:hardware}

NavHAL is designed to be battle-tested on real-world hardware. The current implementation provides first-class support for several popular MCU families.

\subsection{STM32 Series (ARM Cortex-M4)}
The STM32 implementation is the most mature, featuring:
\begin{itemize}
    \item Full register-level abstraction for GPIO, UART, and RCC.
    \item Support for the STM32F4 series, including the ubiquitous F401RE (Nucleo-64).
    \item Hybrid approach: Direct register access for performance-critical paths and LL/HAL wrappers for complex peripherals like USB or Ethernet.
\end{itemize}

\subsection{RP2040 (Raspberry Pi Pico)}
NavHAL provides a lightweight wrapper for the RP2040's unique dual-core architecture and PIO (Programmable I/O) blocks, allowing it to fit seamlessly into the NavHAL ecosystem.

\subsection{AVR ATmega Series}
For 8-bit applications where resource constraints are paramount, NavHAL's AVR implementation focuses on ultra-low memory overhead while maintaining API compatibility with its 32-bit counterparts.

\subsection{Unified Build System: CMake}
NavHAL uses CMake as its primary build system, enabling:
\begin{itemize}
    \item \textbf{Cross-Compilation}: Easy switching between ARM-GCC, AVR-GCC, and other toolchains.
    \item \textbf{Sample-Based Workflow}: A "samples" directory allowing developers to quickly test features (e.g., \texttt{-DSAMPLE=hal\_uart}).
    \item \textbf{Integration}: Seamlessly integrates with CI/CD pipelines (GitHub Actions) for automated testing.
\end{itemize}
